% ============================================================================
% Fichier  : partie5/P05_C29_audit_preuve_investigation.tex
% Titre    : De l'Audit à la Preuve Judiciaire
% Auteur   : MINKA MI NGUIDJOÏ Thierry Emmanuel
% Version  : 1.0 -- novembre 2026
% ============================================================================
% Base : ACSI C20, section « Du constat d'audit à la preuve judiciaire »
% (procédure en cinq étapes ; dossier des ajustements PT-17 remis au
% ministère public, brouillon « fraude » non remis).
% Sources : ISO/IEC 27037:2012 (identification, collecte, acquisition,
% préservation) ; loi n° 2010/012, art. 52 (3), 53 (1) et (3), 55 (1), texte
% officiel lu le 06/11/2026 ; UE SEC4031 et SEC4052 (investigation
% numérique) pour l'approfondissement.
% Corrigés : corriges/Corr_C29.tex.
% ============================================================================

\chapter{De l'Audit à la Preuve Judiciaire}
\label{chap:audit-preuve}
\chaptermark{De l'audit à la preuve}

\epigraph{%
    « L'auditeur constate pour son commanditaire.\\
    Le juge tranche sur des pièces que d'autres ont préservées. »%
}{Principe de la frontière}

\niveauChapitre{Intermédiaire}{%
    Maîtrise des Chapitres~\ref{chap:fraude} (conduite à tenir devant des
    indices), \ref{chap:preuve} (intégrité des extractions) et
    \ref{chap:droit} (constater n'est pas qualifier)%
}

\begin{boxobjectifs}
\begin{itemize}
    \item Distinguer la preuve d'audit et la preuve judiciaire, et leurs
          exigences.
    \item Décrire les quatre étapes de traitement d'une preuve numérique
          selon l'ISO/IEC~27037.
    \item Tenir une chaîne de conservation et un registre des accès.
    \item Connaître les dispositions de procédure de la loi n°~2010/012
          qui concernent l'auditeur.
    \item Remettre un dossier à une autorité judiciaire, et savoir ce
          qu'il ne faut pas remettre.
    \item Situer le rôle de l'expert judiciaire et celui de l'enquêteur.
\end{itemize}
\end{boxobjectifs}

\bigskip

% ============================================================================
\section*{Question directrice : le dossier d'audit tiendra-t-il devant un juge~?}
% ============================================================================

Les constats d'AUDIX CONSEIL sur les ajustements de stock de Douala et les
paiements de TRANSLOG SERVICES ont été transmis au conseil, puis, par le
commissaire aux comptes, au ministère public. Un officier de police
judiciaire demande au cabinet ses pièces. M.~BEYALA sait que son dossier a
été construit pour un audit, pas pour un procès~: il doit savoir ce qu'il
vaut, comment le remettre, et où s'arrête son rôle.

\medskip
\textbf{Plan du chapitre.}
\begin{enumerate}
    \item Preuve d'audit et preuve judiciaire (section~\ref{sec:c29-deux}).
    \item Les quatre étapes de l'ISO/IEC~27037 (section~\ref{sec:c29-27037}).
    \item La chaîne de conservation (section~\ref{sec:c29-chaine}).
    \item Ce que dit la loi n°~2010/012 (section~\ref{sec:c29-loi}).
    \item Remettre un dossier (section~\ref{sec:c29-remettre}).
    \item Auditeur, enquêteur, expert (section~\ref{sec:c29-roles}).
\end{enumerate}

\separateur

% ============================================================================
\section{Preuve d'Audit et Preuve Judiciaire}
\label{sec:c29-deux}
% ============================================================================

\begin{tableaucours}{L{3.0cm}L{4.6cm}L{4.6cm}}{%
    \tableauHeader{} &
    \tableauHeader{Preuve d'audit} &
    \tableauHeader{Preuve judiciaire}
}
Destinataire & Le commanditaire, le comité d'audit & Le juge, les parties
au procès \\
Exigence & Suffisante et appropriée pour une conclusion (Chapitre~\ref{chap:preuve})
& Recevable, intègre, de provenance établie, discutée contradictoirement \\
Contestation & Le contradictoire de la mission & La défense, qui a intérêt à
écarter la pièce \\
Défaut fatal & Une preuve insuffisante affaiblit la conclusion & Une rupture
dans la chaîne de conservation peut faire écarter la pièce \\
\end{tableaucours}
\vspace{-0.8em}
\begin{center}{\small\itshape Tableau~29.1. Preuve d'audit et preuve
judiciaire}\end{center}

Un constat d'audit est opposable s'il peut être défendu devant un tiers
(axiome~6). Une preuve judiciaire doit en outre convaincre un juge que la
pièce n'a pas été altérée depuis sa collecte, face à une partie qui a
intérêt à le contester. Les règles de l'audit (empreintes, lettres
d'extraction, papiers de travail) sont une bonne base~; elles ne suffisent
pas toujours.

% ============================================================================
\section{Les Quatre Étapes de l'ISO/IEC 27037}
\label{sec:c29-27037}
% ============================================================================

L'ISO/IEC~27037 donne les lignes directrices pour traiter une preuve
numérique en quatre étapes~\cite{iso27037_2012}.

\begin{tableaucours}{L{2.6cm}L{5.0cm}L{4.6cm}}{%
    \tableauHeader{Étape} &
    \tableauHeader{Principe} &
    \tableauHeader{Chez AUDIX CONSEIL}
}
Identification & Reconnaître les supports et les données qui peuvent
constituer une preuve & Extraction des ajustements, journaux, courriels \\
Collecte & Prendre les supports physiques, quand c'est possible et
légitime & Non applicable~: l'auditeur ne saisit pas de supports \\
Acquisition & Produire une copie fidèle, vérifiable par empreinte & Lettre
d'extraction, empreinte SHA-256 à la réception \\
Préservation & Protéger la copie contre toute altération, et tracer chaque
accès & Copie de travail, original scellé, registre \\
\end{tableaucours}
\vspace{-0.8em}
\begin{center}{\small\itshape Tableau~29.2. Les quatre étapes de l'ISO/IEC~27037}\end{center}

\begin{boxinfo}
\textbf{Ce que l'auditeur fait déjà.} La lettre d'extraction et
l'empreinte du Chapitre~\ref{chap:preuve} couvrent l'acquisition~; la règle
« travailler sur une copie » couvre une partie de la préservation. Ce qui
manque souvent~: le registre des accès, qui prouve que personne n'a touché
à l'original depuis l'extraction.
\end{boxinfo}

% ============================================================================
\section{La Chaîne de Conservation}
\label{sec:c29-chaine}
% ============================================================================

\begin{boxdef}
\termecle{Chaîne de conservation}\index{chaîne de conservation}~: suite
ininterrompue d'enregistrements qui retrace, pour chaque pièce, qui l'a
détenue, quand, où et pour quoi faire, depuis sa collecte ou son
acquisition jusqu'à sa production en justice.
\end{boxdef}

\begin{boxprocedure}
\textbf{Le registre de conservation d'une pièce}
\begin{itemize}
    \item Identifiant de la pièce, description, source, date et heure
          d'acquisition, empreinte initiale~;
    \item lieu de conservation de l'original et de ses copies~;
    \item chaque accès~: date, heure, personne, motif, empreinte vérifiée
          avant et après~;
    \item chaque remise~: destinataire, qualité, récépissé.
\end{itemize}
\end{boxprocedure}

\begin{boxalert}
\textbf{Violation typique.} Ouvrir l'original d'une extraction « juste pour
vérifier une ligne », sans le consigner.

\smallskip
\textbf{Règle~:} on ne travaille que sur des copies~; l'original n'est
ouvert que pour en vérifier l'empreinte, et chaque ouverture est inscrite au
registre. Une ouverture non tracée est une brèche que la défense
exploitera.
\end{boxalert}

% ============================================================================
\section{Ce que Dit la Loi n° 2010/012}
\label{sec:c29-loi}
% ============================================================================

\begin{boxjuris}
\textbf{\loicam{n°~2010/012}} (texte officiel)

\medskip
\textbf{Article~52, alinéa~3}~: \textit{« Les Officiers de Police Judiciaire
et les agents habilités de l'Agence peuvent, lors des investigations,
accéder aux moyens de transport, à tout local à usage professionnel, à
l'exclusion des domiciles privés, en vue de rechercher, de constater les
infractions, de demander la communication de tous les documents
professionnels et en prendre copie, recueillir, sur convocation ou sur
place, les renseignements et justifications. »}

\medskip
\textbf{Article~53, alinéa~1}~: \textit{« Les perquisitions en matière de
cybercriminalité sont susceptibles de porter sur des données qui peuvent
être des supports physiques ou des copies réalisées en présence des
personnes qui assistent à la perquisition. »}

\medskip
\textbf{Analyse.} La loi confie l'enquête aux officiers de police
judiciaire et aux agents habilités~; elle prévoit qu'ils demandent la
communication des documents professionnels et en prennent copie. Un
cabinet d'audit qui reçoit une telle demande y répond dans le cadre fixé
par la procédure (forme de la demande, secret professionnel, conseil
juridique~: Chapitre~\ref{chap:droit}). L'article~55 permet en outre aux
autorités judiciaires de requérir une personne qualifiée pour rendre
lisibles des données chiffrées~: c'est une mission d'expertise, distincte
de l'audit.
\end{boxjuris}

% ============================================================================
\section{Remettre un Dossier}
\label{sec:c29-remettre}
% ============================================================================

L'auditeur n'enquête pas~: il préserve ce qu'il a, et le remet. Un
auditeur qui « complète » le dossier en interrogeant des suspects ou en
extrayant de nouvelles données hors de sa mission fragilise la procédure.

\begin{boxprocedure}
\textbf{Remettre un constat à l'autorité judiciaire}
\begin{enumerate}
    \item \textbf{Geler}~: ne plus modifier les fichiers concernés~; en
          faire une copie de travail, l'original restant intact.
    \item \textbf{Prouver l'intégrité}~: empreinte de chaque fichier,
          vérifiée et consignée (Chapitre~\ref{chap:preuve}).
    \item \textbf{Documenter la provenance}~: lettre d'extraction (source,
          date, requête, paramètres), papiers de travail, observation.
    \item \textbf{Tenir le registre}~: chaque remise, chaque accès, avec la
          date, la personne et l'empreinte vérifiée.
    \item \textbf{Remettre}~: à l'autorité qui le demande, contre récépissé,
          sans commentaire qualifiant les faits (Chapitre~\ref{chap:droit}).
\end{enumerate}
\end{boxprocedure}

\begin{boxscenario}{Opération LIZBIZ. Le dossier des ajustements}
Le ministère public demande les pièces relatives aux ajustements de stock
de Douala. AUDIX CONSEIL remet l'extraction PT-17 et son empreinte, la
lettre d'extraction, le papier de travail du test et l'observation, avec
un registre qui retrace chaque accès depuis l'extraction. Il ne remet pas
le brouillon de l'auditeur junior qui parlait de « fraude »~: ce n'est
pas une pièce de la mission, et il qualifiait les faits sans en avoir
le droit.
\end{boxscenario}

\begin{boxinfo}
\textbf{Et si l'autorité demande tout~?} La remise relève d'une décision
juridique, prise avec le conseil du cabinet~: l'auditeur ne choisit pas
seul ce qu'il retient. Le principe du cours est de documenter précisément
ce qui est remis, ce qui ne l'est pas et pourquoi, pour que la décision
puisse être revue.
\end{boxinfo}

% ============================================================================
\section{Auditeur, Enquêteur, Expert}
\label{sec:c29-roles}
% ============================================================================

\begin{tableaucours}{L{2.6cm}L{3.6cm}L{3.0cm}L{3.0cm}}{%
    \tableauHeader{Rôle} &
    \tableauHeader{Mission} &
    \tableauHeader{Mandant} &
    \tableauHeader{Pouvoirs}
}
Auditeur & Évaluer et conclure sur un objet, selon des critères & Le
commanditaire de la mission & Ceux de la lettre de mission \\
Enquêteur & Rechercher et constater les infractions & L'autorité
judiciaire & Ceux du code de procédure pénale et de la loi
n°~2010/012 \\
Expert judiciaire & Répondre aux questions techniques du juge & Le juge &
Ceux que lui donne la décision qui le désigne \\
\end{tableaucours}
\vspace{-0.8em}
\begin{center}{\small\itshape Tableau~29.3. Trois rôles distincts}\end{center}

Une même personne peut exercer ces rôles à des moments différents, jamais
dans la même affaire sans que la procédure le permette~: un auditeur qui
deviendrait l'expert judiciaire d'un dossier qu'il a audité poserait un
problème d'indépendance (Chapitre~\ref{chap:deontologie}). L'investigation
numérique proprement dite (acquisition de supports, analyse de mémoire,
reconstitution de chronologies) fait l'objet des UE~SEC4031 et SEC4052.

\begin{boxmetier}{Auditeur SI}
\textbf{Construire un dossier qui pourra servir}
\begin{enumerate}
    \item Calculer l'empreinte de toute extraction à la réception, et la
          noter dans la lettre d'extraction.
    \item Travailler sur des copies~; sceller l'original.
    \item Tenir un registre des accès dès qu'un constat touche à une
          irrégularité possible.
    \item Séparer constatation et analyse, sans qualification pénale.
    \item Se retirer de l'enquête~: remettre, ne pas enquêter.
\end{enumerate}
\end{boxmetier}

\begin{boxmetier}{Responsable de la conformité}
\textbf{Organiser la réponse aux autorités.} Le responsable de la
conformité désigne qui reçoit les demandes des autorités, qui décide de la
remise (avec le conseil juridique), et comment elle est tracée. Une
demande reçue par un agent isolé, sans procédure, est un risque pour
l'organisation et pour la preuve.
\end{boxmetier}

\begin{boxinfo}
\textbf{Ce que dit la recherche.} L'axiome d'opposabilité d'A0 trouve ici
son épreuve la plus dure~: une partie qui a intérêt à écarter la pièce, et
un juge qui en apprécie la valeur~\cite{minka2026auditologie}. La
discipline de l'audit (empreinte, provenance, séparation des faits et de
l'analyse) prépare à cette épreuve, sans dispenser des règles propres à la
procédure.
\end{boxinfo}

\separateur

\begin{boxresume}
\textbf{Ce qu'il faut retenir de ce chapitre~:}
\begin{itemize}
    \item \textbf{La preuve judiciaire} ajoute à la preuve d'audit la
          recevabilité, la provenance établie et la discussion par une
          partie adverse.
    \item \textbf{L'ISO/IEC~27037}~: identification, collecte, acquisition,
          préservation.
    \item \textbf{La chaîne de conservation} trace chaque accès et chaque
          remise~; une brèche peut faire écarter la pièce.
    \item \textbf{La loi n°~2010/012} confie l'enquête aux officiers de
          police judiciaire et aux agents habilités~; l'auditeur répond dans
          le cadre de la procédure, avec un conseil juridique.
    \item \textbf{L'auditeur préserve et remet, il n'enquête pas}~;
          l'expert judiciaire répond au juge, dans une autre mission.
\end{itemize}
\end{boxresume}

\begin{boxinfo}
\textbf{Points de contrôle de l'élève.} Je sais~: distinguer preuve d'audit
et preuve judiciaire~; nommer les quatre étapes de l'ISO/IEC~27037~; tenir
un registre de conservation~; citer les articles~52 et~53~; remettre un
dossier en cinq étapes~; distinguer auditeur, enquêteur et expert.
\end{boxinfo}

\bigskip

\begin{boxexo}[title={\ding{45}\quad Exercice~29.1. Audit ou justice~? \niveaubase}]
Pour chaque pièce, dites si elle suffit pour une conclusion d'audit, et ce
qu'il lui manque pour une procédure judiciaire~: (a) une extraction reçue
par courriel, avec une empreinte calculée trois jours plus tard~; (b) une
capture d'écran d'un journal~; (c) une extraction faite devant l'auditeur,
avec lettre d'extraction et empreinte immédiate.
\end{boxexo}

\begin{boxexo}[title={\ding{45}\quad Exercice~29.2. Les étapes \niveaubase}]
Rangez dans les étapes de l'ISO/IEC~27037~: (a) calculer l'empreinte d'une
copie~; (b) repérer qu'un serveur de messagerie contient des échanges
utiles~; (c) placer l'original sur un support scellé~; (d) emporter un
ordinateur portable sous scellé.
\end{boxexo}

\begin{boxexo}[title={\ding{45}\quad Exercice~29.3. Le registre \niveaumoyen}]
Rédigez les cinq premières lignes du registre de conservation de
l'extraction PT-17, de son acquisition à sa remise au ministère public.
\end{boxexo}

\begin{boxexo}[title={\ding{45}\quad Exercice~29.4. Une brèche \niveaumoyen}]
L'auditeur junior a ouvert l'original de PT-17 dans un tableur pour
« vérifier une ligne », puis l'a refermé sans l'enregistrer. Le registre
n'en dit rien. Quelles conséquences, et que faire maintenant~?
\end{boxexo}

\begin{boxexo}[title={\ding{45}\quad Exercice~29.5. Ce qu'on ne remet pas \niveaumoyen}]
Justifiez, pièce par pièce, la remise ou non~: le papier de travail du
test~; le brouillon qui parle de « fraude »~; les courriels internes de
l'équipe~; l'observation après contradictoire~; la note au commanditaire du
12~juin.
\end{boxexo}

\begin{boxexo}[title={\ding{45}\quad Exercice~29.6. Article 52 \niveaumoyen}]
Un agent se présente au cabinet et demande « tous les documents » sur
LiZbiZ, sans écrit. Que vérifie l'associé avant de répondre, au regard de
l'article~52 et de la procédure~? Que consigne-t-il~?
\end{boxexo}

\begin{boxexo}[title={\ding{45}\quad Exercice~29.7. Trois rôles \niveaubase}]
M.~BEYALA est sollicité pour être expert judiciaire dans l'affaire LiZbiZ.
Peut-il accepter~? Justifiez par le Tableau~29.3 et le
Chapitre~\ref{chap:deontologie}.
\end{boxexo}

\begin{boxexo}[title={\ding{45}\quad Exercice~29.8. Un dossier qui servira \niveaumoyen}]
Pour l'audit de la BANQUE SANAGA, Mme~EKAMBI découvre qu'un agent a modifié
un virement après sa validation. Que fait-elle, dès cet instant, pour que
ses pièces puissent servir si une plainte est déposée~?
\end{boxexo}

\begin{boxexo}[title={\ding{45}\quad Exercice~29.9. La défense attaque \niveauexpert}]
L'avocat de la défense soutient que l'extraction PT-17 a pu être modifiée,
puisqu'elle a été produite par l'entreprise elle-même. Rédigez la réponse
que permettent les pièces du dossier d'AUDIX CONSEIL, et dites ce qui la
renforcerait.
\end{boxexo}

\begin{boxexo}[title={\ding{45}\quad Exercice~29.10. Une procédure interne \niveauexpert}]
Rédigez pour AUDIX CONSEIL une procédure interne d'une page~: « Que faire
quand une mission révèle des indices d'infraction~? », de la détection à la
remise, en distinguant les rôles de l'auditeur, du chef de mission, de
l'associé et du conseil juridique.
\end{boxexo}

\bigskip

\begin{boxinfo}
\textbf{Pour aller plus loin.}
\begin{itemize}
    \item ISO/IEC~27037:2012~\cite{iso27037_2012}.
    \item Loi n°~2010/012~\cite{loi2010012}, articles~52 à~59.
    \item Les UE~SEC4031 (Introduction à l'investigation numérique) et
          SEC4052 (Investigation numérique avancée) de l'ENSPY.
\end{itemize}
\end{boxinfo}

\begin{boxinfo}
\textbf{Transition.} La Partie~\ref{part:avance} a montré les défis de
l'audit d'aujourd'hui. La Partie~\ref{part:recherche} revient aux
fondements~: le programme de recherche de l'auditologie, puis la mesure de
l'incertitude et de la stratégie.
\end{boxinfo}
